\chapter{Comparison methodology}
\label{ch:methodology}

This chapter is written to be difficult to misquote. It states what is measured,
what each measurement licenses, and what it does not.

\section{Why the naive comparison is misleading}

The tempting experiment is to run a solver on the CPU, run it on the GPU, divide
the times and print the ratio. Three things go wrong.

\paragraph{It compares different algorithms.} Natural ordering Gauss Seidel
cannot run on a GPU at all. If the CPU runs Gauss Seidel and the GPU runs
Jacobi, the ratio conflates the device with the fact that Jacobi needs about
twice as many iterations to reach the same tolerance.

\paragraph{It compares different amounts of effort.} A carefully written GPU port
measured against a CPU version compiled without vectorisation or threading
produces an impressive number about nothing.

\paragraph{It hides what actually limits the calculation.} A five point stencil
sweep moves three doubles and performs about six floating point operations per
unknown. It is bandwidth bound on every machine ever built, so a ratio of
runtimes is, to within a small correction, a ratio of memory bandwidths dressed
up as a statement about parallel programming.

\section{The protocol}

\subsection{First: iterations to tolerance}

Iteration counts are properties of the mathematics. They are identical on every
machine that runs them, they are what determine whether one method beats another
in principle, and they are therefore reported before any timing. All methods stop
on the same criterion, a relative residual below $10^{-8}$, evaluated the same
way by the same shared driver, so no method is flattered by a laxer test.

\subsection{Like against like}

GPU Jacobi is compared against CPU Jacobi. GPU red black Gauss Seidel is
compared against CPU red black Gauss Seidel. The cost of the red black
reordering, which is a change of iteration matrix and therefore of iteration
count, is measured separately and charged to both sides equally rather than
being allowed to inflate the device figure.

\subsection{Second: bytes per unknown, counted}

A Jacobi sweep writes one value, reads one right hand side entry and reads the
previous iterate. The four stencil neighbours of consecutive unknowns overlap,
so in the streaming limit each is already resident from a nearby access. That is
three doubles, \textbf{24 bytes per unknown per pass}. The same figure is
returned by the host implementation and set by the device solver, so the number
used in the analysis comes from the code rather than from an estimate.

That count is the conservative one and it is not the only defensible one. A
store to a line the cache does not hold has to fetch that line before it can
modify it, which is read for ownership and is a property of a write allocate
cache rather than of the algorithm. A Jacobi pass writes an output array it
never reads, so under that model the count is short by one double per unknown
and the true figure is \textbf{32 bytes}. Which of the two is right is an open
question, so both are published: every result row carries both counts, and every
table and figure here that shows an achieved bandwidth shows both, labelled
\emph{declared} and \emph{counted}. Neither replaces the other, and quoting one
without the other is quoting half the evidence.

\subsection{Pre registered: the rule that settles the traffic model}
\label{sec:preregistered-traffic}

The question is settled by measurement and the rule that reads that measurement
was fixed before it was taken, in
\texttt{benchmarks/sweep\_matrix.yaml} under
\texttt{preregistered.traffic\_model}. Two host STREAM triads run over the same
arrays, the same sizes, the same worker counts and the same repetitions,
differing in the store instruction and in nothing else: a plain C++ loop, whose
compiler emitted stores are the ones under suspicion, and the same triad written
through \texttt{\_mm256\_stream\_pd} with one \texttt{\_mm\_sfence}, which does
not fetch the line it overwrites. Both report against the same declared 24 bytes
per element, so the ratio of the two is the ratio of the traffic they really
move.

The rule, quoted from the registration:

\begin{quote}
A ratio above 1.20 selects the read for ownership model, 32 bytes per unknown
per pass. A ratio below 1.10 selects the conservative model, 24 bytes. A ratio
from 1.10 to 1.20 inclusive is recorded as unresolved and both models are
carried. The thresholds are fixed here, before any measurement, and are not
revisited afterwards.
\end{quote}

The statistic the rule is applied to was amended on 2026-09-06, before the
publication session and before any measurement the rule governs, because the
statistic as first registered had a run to run spread as wide as the undecided
band it had to fall inside, and because its two halves could be taken at
different worker counts. The thresholds and the sentences below were not
touched. As amended: five repetitions of each probe at every worker count,
alternating the two so a repetition's arms see the same machine; $w^*$ is the
worker count at which the plain triad's median is highest; the statistic is the
median over the five repetitions of the non temporal figure over the plain
figure at $w^*$, each ratio taken within one repetition; and if the interval
from the smallest to the largest of those five ratios contains either threshold
the outcome is unresolved whatever the statistic is, which is ground rule 7
applied to a denominator.

\paragraph{Pre registered, one sentence per outcome.} Each of the following was
written before the measurement, and the report carries whichever one the
measurement selects, with the ratio substituted for the placeholder. They are
quoted here verbatim so that a reader can check the report against what was
promised rather than against what was convenient.

\begin{quote}
\textbf{If the read for ownership model is selected.} The non temporal triad
reached a ratio of \texttt{<ratio>} against the plain triad on the publication
machine, above the 1.20 threshold fixed before the measurement, so the read for
ownership model is selected: a Jacobi pass over the five point stencil moves 32
bytes per unknown and not 24, the counted column is the achieved bandwidth this
report compares against each device's own triad, and the host figures rise by a
third while the device figures do not move, which is why host and device
efficiency converge on the Jacobi row.
\end{quote}

\begin{quote}
\textbf{If the conservative model is selected.} The non temporal triad reached a
ratio of \texttt{<ratio>} against the plain triad on the publication machine,
below the 1.10 threshold fixed before the measurement, so the conservative model
is selected: a Jacobi pass moves 24 bytes per unknown, the declared column
stands as the achieved bandwidth of this report, and Section 4.2's read for
ownership argument is refuted on this machine rather than confirmed. The Jacobi
efficiency gap between host and device is then a real gap and not an artefact of
the denominator, and the only correction this release makes to the work unit is
the pass count of phase A2.
\end{quote}

\begin{quote}
\textbf{If it is unresolved.} The non temporal triad reached a ratio of
\texttt{<ratio>} against the plain triad on the publication machine, between the
1.10 and 1.20 thresholds fixed before the measurement, so the traffic model is
recorded as unresolved: both counts, 24 and 32 bytes per unknown per pass, are
carried in every table and figure of this report, neither is presented on its
own as the achieved bandwidth, and no claim is made that rests on one of them
and would fail under the other.
\end{quote}

Which of the three the measurement selected is stated in Chapter
\ref{ch:results}, from the session manifest, by the generator rather than by a
hand that had already seen the number. The publication session applies the rule
for this release; the assembly triad of release 1.2.0 rebuilds the same
experiment and confirms the outcome rather than gating it.

\subsection{Third: each device's own bandwidth, measured}

Both devices are probed with the same STREAM triad kernel,
$a_i = b_i + q\,c_i$, which moves two reads and one write per element and has no
reuse, so it measures the memory system and nothing else \cite{mccalpin1995}.
The host probe runs over the execution backend rather than on one thread,
because a single core cannot saturate a twenty core part and dividing by a
single core figure would make every parallel result look superlinear.

Manufacturer bandwidth figures are never used. They are a bus width times a
clock, no real kernel reaches them, and dividing by one would make both devices
look equally inefficient while saying nothing about which is being used well.

\subsection{The comparable number}

\begin{equation}
  \text{efficiency} =
  \frac{\text{achieved GiB/s on the sweep}}
       {\text{that device's own measured triad GiB/s}} .
  \label{eq:efficiency}
\end{equation}

This is the figure to compare across devices. It is dimensionless, bounded above
by one, and it transfers: a reader with different hardware can measure their own
triad, apply the same fraction, and predict what they would get. It is the
standard roofline reasoning \cite{williams2009} applied to a kernel whose
arithmetic intensity is fixed and known.

\subsection{Last: absolute time}

Seconds are reported last and labelled as a property of one specific CPU paired
with one specific GPU, under one set of compiler flags, on one day. The measured
speedup is then \emph{decomposed} into the factor attributable to the memory
system and the factor attributable to the implementation, so that a reader can
see how much of it is the graphics card and how much of it is the port.

\section{Measuring at sizes where methods do not converge}

Section 8.3 of the specification fixes the comparison at 1024, 2048 and 4096
squared unknowns. At those sizes the stationary methods cannot be run to
convergence: Jacobi on a 1024 squared grid needs of order four million sweeps,
which is days of compute for a number the closed form already predicts.

The two quantities are therefore measured separately, which is more rigorous
than measuring them together rather than less:

\begin{itemize}
  \item \textbf{Iteration counts} are measured at 63, 127 and 255, where running
        to convergence is affordable, and checked against the closed form rates
        of Chapter \ref{ch:background}. It is that agreement which makes
        extrapolating them legitimate rather than a guess. The methods whose
        counts grow like $n$ rather than $n^2$, namely optimally relaxed SOR, its
        red black form and conjugate gradient, are additionally measured at 511
        and 1023, which demonstrates the change of growth order directly.
  \item \textbf{Per iteration cost} is measured at the large sizes with a fixed,
        stated number of sweeps. This is what the bandwidth analysis needs, and
        it is exactly what is affordable there.
\end{itemize}

Time to solution then follows from the two, and because both are tabulated the
crossover point between any two methods on any two devices can be \emph{derived}
rather than asserted:

\begin{equation}
  t = \frac{\text{iterations}(\text{method}, n)\;\times\;n^2\;\times\;\text{bytes per unknown}}
           {\text{achieved bandwidth}(\text{device}, \text{method}, n)} .
  \label{eq:crossover}
\end{equation}

\section{Experimental hygiene}

\begin{itemize}
  \item Every configuration is run once untimed as a warm up, so page faults and
        first touch placement do not land in the measurement, then repeated five
        times with the median reported and the minimum and maximum recorded.
  \item Every result row carries the problem, solver, backend, worker count,
        rank count, pinning policy, reduction mode, schedule, seed and commit
        hash. A row cannot claim provenance the binary did not have: the commit
        is stamped into the binary at configure time and a dirty working tree is
        recorded as such.
  \item Problems are generated from recorded seeds, so any row can be rebuilt.
  \item Both bandwidth probes run once per sweep session and land in the session
        manifest alongside the toolchain versions.
  \item A configuration that fails, or that does not apply, is recorded as such
        with its message rather than dropped.
  \item Compiled with \texttt{-O3 -march=native} and without fast math. Fast math
        would licence the compiler to reassociate the reductions that the
        equivalence suite exists to check.
\end{itemize}

\section{Comparing these numbers with release 1.0.0's}
\label{sec:across-releases}

Every timing in this report replaces the corresponding timing of release 1.0.0
rather than extending it, and none of the values does.

\paragraph{The values are bit identical.} Release 1.1.0 keeps
\texttt{PNL\_REDUCTION\_ACCUMULATORS} at one, so the reduction is the same
ascending chain it was, and every iterate and every residual this library
computes matches release 1.0.0's bit for bit. That is a supported configuration
and not a leftover. The reduction association changes only when that number
changes, which is a later release, and its changelog entry leads with the escape
hatch rather than burying it.

\paragraph{Every timing moved, for three deliberate reasons.} The publication
compiler changed from an unreleased trunk snapshot to a released one, and a
change of compiler changes every timing. The Jacobi solver no longer copies its
whole state on the calling thread once per iteration, so its rows measure less
work rather than the same work faster. The solver's workspace is allocated
outside the timed region, so a repetition no longer pays first touch page fault
cost inside the measurement. The work unit moved with them: sweeps and passes
are separate quantities now, Richardson is charged for one residual evaluation
rather than two, and the symmetric methods are charged the two sweeps they
perform. A 1.1.0 row is therefore not the same quantity as a 1.0.0 row even
where the seconds are close.

The before and after for every headline number that moved is recorded phase by
phase in \texttt{PROGRESS.md}, against the phase that moved it and beside the
gate output that measured it. It is not copied into this chapter, because a
number typed here stops being true at the next sweep, which is the fault this
chapter exists to prevent. The superseded generation keeps its own summary and
its own manifest under \texttt{experiments/results/archive/}, so either can be
rebuilt.

\section{What this comparison does not establish}

\begin{itemize}
  \item \textbf{Nothing about single solve latency.} Transfer time is reported
        separately and is not folded into the bandwidth figure. For one small
        solve, moving the problem across PCIe can cost more than the solve; for a
        simulation that keeps state resident it is paid once.
  \item \textbf{Nothing about other problem classes.} These conclusions hold for
        bandwidth bound stencil work. A dense factorisation is compute bound and
        would rank the devices differently.
  \item \textbf{Nothing about other hardware.} Both measured bandwidths are
        published so that a reader can substitute their own.
  \item \textbf{Nothing about power, cost or programmer time.} Not measured, so
        not claimed.
  \item \textbf{Nothing about performance from the continuous integration
        badge.} CI runs on a shared virtual machine with no GPU; it checks
        correctness, and every timing in this report comes from the target
        machine.
\end{itemize}
